%%
%% This is a placeholder manuscript shell for ACL 2027 (via ACL Rolling Review - ARR Jan 2027).
%% Refresh this manuscript before submission by syncing with the research repository upon activation.
%%
\documentclass[11pt]{article}
\usepackage[review]{acl}
\usepackage{times}
\usepackage{latexsym}
\usepackage{microtype}
\usepackage{booktabs}
\usepackage{amsmath}

\title{Prompt Injection and Multilingual Representation Disparity: Calibration Asymmetry in Cross-Lingual Safety Evaluation}

\author{Anonymous ACL Submission}

\begin{document}
\maketitle

\begin{abstract}
This study investigates cross-lingual vulnerability and detection dynamics for prompt injection in language models and autonomous agents when input contexts span typologically diverse languages. Evaluating 1,160 InjecAgent episodes under locally served open-weight models reveals near-zero attack success rates across all languages, demonstrating an experimental floor effect rather than genuine cross-lingual robustness. Across 57,912 scored inferences evaluating open injection detectors, cross-lingual recall collapse does not occur at English-calibrated decision thresholds. Instead, multilingual encoders display severe calibration asymmetry: raw injection scores are systematically elevated on non-English text, yielding benign false-alarm rates up to 67.7\% in Chinese and 15.7\% in Arabic. Applying language-conditioned thresholds causes Arabic detection recall to plunge from 54.0\% to 7.0\%, revealing an acute trade-off between false-alarm suppression and evasion resistance. Tokenisation fertility analysis indicates up to 4.33x token expansion in non-Latin scripts, highlighting subword fragmentation as a primary correlate of representation drift. These empirical findings indicate that multilingual safety evaluation requires language-specific calibration and script-aware representation auditing.
\end{abstract}

\section{Introduction}
As large language models and autonomous agentic systems are deployed globally, evaluating safety boundaries across diverse languages is essential. Prior research posits that safety alignments established predominantly on English corpora may suffer from cross-lingual degradation when processing non-English inputs. This work empirically evaluates whether prompt injection attacks exhibit cross-lingual transfer disparities, analysing both model execution and detection dynamics.

\section{Methodology and Experimental Framework}
We formulate a controlled cross-lingual injection benchmark evaluating two complementary dimensions:
\begin{enumerate}
    \item \textbf{Agent Execution:} 1,160 episodes derived from InjecAgent across English, Chinese, Arabic, and Spanish using locally served open-weight models.
    \item \textbf{Detector Calibration:} 57,912 scored inferences evaluating open-source injection detectors under varying language-specific thresholding regimes.
\end{enumerate}

\section{Empirical Results}
Our evaluation yields two key findings:
\begin{itemize}
    \item \textbf{Execution Floor Effect:} Attack success remains at or near zero across all evaluated languages (English 1.7\% and 0.0\%), precluding claims of an execution gap.
    \item \textbf{Calibration Asymmetry:} Detectors do not experience cross-lingual recall collapse; rather, multilingual encoders exhibit elevated anomaly scores on non-English text, producing false-positive rates of 67.7\% on benign Chinese spans.
\end{itemize}

\section{Tokenisation Fertility and Script Disparities}
We examine subword tokenisation fertility across scripts, observing up to 4.33x token expansion in non-Latin scripts (e.g. Arabic and Chinese) relative to English. We analyse this syntactic fragmentation as a structural correlate of calibration disparity.

\section{Limitations}
This evaluation is bounded by the open-weight models evaluated and synthetic benchmark episodes. The floor effect observed in agent execution highlights the need for more challenging agentic environments before definitive conclusions regarding operational execution disparities can be drawn.

\section{Ethics Statement}
All injection payloads evaluated in this work are derived from defensive research benchmarks and synthetic teaching fixtures. No proprietary models or private user data were accessed. This research aims to enhance the security and fairness of multilingual language model deployments.

\bibliography{custom}

\end{document}
